\documentclass[12pt]{article}

\usepackage[a4paper,margin=1in]{geometry}
\usepackage{amsmath,amssymb}
\usepackage{setspace}
\usepackage{titlesec}

\setstretch{1.2}

\title{CSE400 – Fundamentals of Probability in Computing\\
Lecture 4: Joint Probability and Conditional Probability}
\author{Instructor: Dhaval Patel, PhD}
\date{Date: January 15, 2026}

\begin{document}

\maketitle

\section*{Lecture - 4}

\section*{1. Engineering Applications of Probability}

Examples of engineering systems where probability is applied include:

\textbf{Speech Recognition System}

A vocabulary such as Hello, Yes, No, Bye can be represented using waveform templates. Different speakers (male, female, child) produce different waveform patterns for the same word.

\textbf{Radar System}

Radar signals are transmitted and reflected from objects such as aircraft. The received signal is processed to determine object detection.

\textbf{Communication Network}

Data is transmitted from a source to a destination through multiple network nodes and devices.

\section*{2. Introduction to Probability Theory}

\subsection*{2.1 Experiments, Sample Space, and Events}

\textbf{Definition: Experiment (E)}

An experiment is a procedure that produces some result.

Example:

Tossing a coin five times $E_5$.

\textbf{Definition: Outcome ($\xi$)}

An outcome is a possible result of an experiment.

Example:

One outcome of $E_5$:

$\xi_1 = HHTHT$

\textbf{Definition: Event}

An event is a set of outcomes of an experiment.

Example:

Let event $C$ in experiment $E_5$ be:

$C = \{$all outcomes consisting of an even number of heads$\}$

\textbf{Definition: Sample Space (S)}

The sample space is the collection or set of all possible distinct outcomes of an experiment.

These outcomes are:

\textbf{Mutually Exclusive}

Example: In a coin flip, you can obtain heads or tails, but not both.

\textbf{Collectively Exhaustive}

All possible outcomes are included in the sample space.

Example: Flip a coin.

\subsection*{Properties of Sample Space}

$S$ is the universal set of outcomes of an experiment.

The sample space can be:

Discrete

Countably infinite

Continuous

\subsection*{Examples of Experiments}

Flipping a fair coin once

Rolling a cubical die with numbered faces

Rolling two dice

Flipping a coin until a tails occurs

Random number generator with interval $[0,1]$

\section*{3. Axioms of Probability}

\textbf{Definition: Probability}

Probability is a measure of the likelihood of events or a function of an event that produces a numerical quantity measuring the likelihood of that event.

\textbf{Axioms}

The following statements are taken as self-evident and require no proof.

For any event $A$:

\[
0 \le Pr(A) \le 1
\]

If $S$ is the sample space:

\[
Pr(S) = 1
\]

If $A$ and $B$ are mutually exclusive $(A \cap B = \emptyset)$:

\[
Pr(A \cup B) = Pr(A) + Pr(B)
\]

For mutually exclusive sets $A_i$:

\[
Pr\left(\bigcup_{i=1}^{\infty} A_i \right) =
\sum_{i=1}^{\infty} Pr(A_i)
\]

\section*{4. Corollary from Probability Axioms}

Let $A$ be the null event for all values of $i$ greater than $n$.

\textbf{Corollary 2.1}

For $M$ mutually exclusive events:

\[
Pr\left(\bigcup_{i=1}^{M} A_i \right) =
\sum_{i=1}^{M} Pr(A_i)
\]

Axiom 3 is equivalent to Corollary 2.1 when the sample space is finite.

The generality of Axiom 3 is necessary when the sample space contains infinitely many points.

\section*{5. Propositions from Probability Axioms}

\textbf{Proposition 2.1}

\[
Pr(A^c) = 1 - Pr(A)
\]

\textbf{Proposition 2.2}

If $A \subset B$:

\[
Pr(A) \le Pr(B)
\]

\textbf{Proposition 2.3}

For any sets $A$ and $B$:

\[
Pr(A \cup B) = Pr(A) + Pr(B) - Pr(A \cap B)
\]

\textbf{Proposition 2.4 (Inclusion–Exclusion Principle)}

\[
Pr(A_1 \cup A_2 \cup ... \cup A_M)
=
\sum_{i=1}^{M} Pr(A_i)
-
\sum_{i_1<i_2} Pr(A_{i_1}A_{i_2})
+ ... +
(-1)^{M+1} Pr(A_1A_2...A_M)
\]

\section*{6. Assigning Probabilities}

\textbf{Classical Approach}

Example 2.6

Coin flipping experiment:

Compute

$Pr(H), Pr(T)$

Example 2.7

Dice rolling experiment:

Compute

$Pr(even\ number\ is\ rolled)$

Example 2.8

Two dice are rolled.

Let event $A$:

Sum of the two dice equals five.

Compute:

$Pr(A)$

\textbf{Relative Frequency Approach}

Example: Dice experiment where event $A$ = sum of two dice equals five.

Simulation results:

\begin{center}
\begin{tabular}{c|cccccccccc}
n & 1000 & 2000 & 3000 & 4000 & 5000 & 6000 & 7000 & 8000 & 9000 & 10000 \\
\hline
$n_A$ & 96 & 200 & 314 & 408 & 521 & 630 & 751 & 859 & 970 & 1095 \\
$n_A/n$ & 0.096 & 0.100 & 0.105 & 0.102 & 0.104 & 0.105 & 0.107 & 0.107 & 0.108 & 0.110
\end{tabular}
\end{center}

\textbf{Inference}

To obtain the exact probability of an event, the experiment must be repeated infinitely many times.

Many random phenomena of interest are not repeatable.

\section*{7. Joint Probability}

\textbf{Motivation}

Events are not always mutually exclusive.

We are interested in the probability of the intersection of events $A$ and $B$.

\textbf{Notation}

Joint probability is denoted as:

\[
Pr(A,B) \quad \text{or} \quad Pr(A \cap B)
\]

Multiple Events

For events $A_1, A_2, ..., A_M$:

\[
Pr(A_1, A_2, ..., A_M)
\]

\section*{8. Example 1: Card Deck}

Consider a deck of 52 playing cards.

Define events:

$A$: Red card selected

$B$: Number card selected (Ace counted as number)

$C$: Heart card selected

Find:

$Pr(A), Pr(B), Pr(C), Pr(A,B), Pr(A,C), Pr(B,C)$

Solution

Individual Probabilities

\[
Pr(A) = \frac{26}{52} = \frac{1}{2}
\]

\[
Pr(B) = \frac{40}{52} = \frac{5}{26}
\]

\[
Pr(C) = \frac{13}{52} = \frac{1}{4}
\]

Joint Probabilities

\[
Pr(A,B) = \frac{20}{52} = \frac{5}{13}
\]

\[
Pr(A,C) = \frac{13}{52} = \frac{1}{4}
\]

\[
Pr(B,C) = \frac{10}{52}
\]

\section*{9. Example 2: Costume Party}

Alex has the following clothing:

Tops:

3 white T-shirts

1 radioactive green cape

Bottoms:

2 flannel dinosaur pajama pants

4 polka-dot boxer shorts

Step 1: Analyze Tops

Total tops:

\[
3+1=4
\]

\[
Pr(Cape)=\frac{1}{4}
\]

Step 2: Analyze Bottoms

Total bottoms:

\[
2+4=6
\]

\[
Pr(Boxers)=\frac{4}{6}=\frac{2}{3}
\]

Step 3: Joint Probability

\[
Pr(Cape \cap Boxers)=Pr(Cape)\times Pr(Boxers)
\]

\[
=\frac{1}{4}\times\frac{4}{6}=\frac{1}{6}
\]

Correct Option: C

\section*{10. Conditional Probability}

\textbf{Motivation}

The occurrence of one event may depend on another event.

\textbf{Definition}

The probability of event $A$ given event $B$:

\[
Pr(A|B)=\frac{Pr(A,B)}{Pr(B)}
\]

where

\[
Pr(B)>0
\]

\textbf{Product Rule}

\[
Pr(A,B)=Pr(A|B)Pr(B)=Pr(B|A)Pr(A)
\]

\section*{11. Example 3: Cards Without Replacement}

Two cards are selected without replacement.

Define:

$A$: First card is a spade

$B$: Second card is a spade

Find:

$Pr(B|A)$

Initial State

Total cards = 52

Spades = 13

After Event $A$

One spade removed.

Remaining:

Total cards = $52-1=51$

Spades = $13-1=12$

Calculation

\[
Pr(B|A)=\frac{12}{51}
\]

\section*{12. Example 4: Game of Poker}

Five cards are dealt.

A flush means all cards are from the same suit.

Part 1: Flush in Spades

\[
P(1st\ spade)=\frac{13}{52}
\]

\[
P(2nd\ spade)=\frac{12}{51}
\]

\[
P(3rd\ spade)=\frac{11}{50}
\]

\[
P(4th\ spade)=\frac{10}{49}
\]

\[
P(5th\ spade)=\frac{9}{48}
\]

\[
P(Spade\ Flush)=
\frac{13}{52}\cdot
\frac{12}{51}\cdot
\frac{11}{50}\cdot
\frac{10}{49}\cdot
\frac{9}{48}
\]

Part 2: Any Flush

There are four suits:

Spades, Hearts, Diamonds, Clubs.

Since events are mutually exclusive:

\[
P(Any\ Flush)=4\times P(Spade\ Flush)
\]

\section*{13. Example 5: The Missing Key}

Bob is 80\% certain the key is in his jacket.

Left pocket:

\[
P(L)=0.4
\]

Right pocket:

\[
P(R)=0.4
\]

In jacket:

\[
P(K)=0.8
\]

Goal:

\[
P(R|L^c)
\]

Probability that the key is in the right pocket given it is not in the left pocket.

Calculation

\[
P(R|L^c)=\frac{P(R \cap L^c)}{P(L^c)}
\]

Since $R$ implies $L^c$:

\[
P(R \cap L^c)=P(R)
\]

\[
P(L^c)=1-P(L)
\]

Therefore

\[
P(R|L^c)=\frac{0.4}{1-0.4}=\frac{0.4}{0.6}=\frac{2}{3}
\]

Correct Option: C

\end{document}
```
